% chuong1.tex — sửa nội dung chương ở đây
\section{Phạm vi và bối cảnh}\label{sec:1}

\subsection{Bối cảnh và giới thiệu tình huống}\label{sec:1.1}

\subsubsection{Dữ liệu cá nhân đã trở thành một loại hàng hóa}\label{sec:1.1.1}

Dữ liệu cá nhân là thông tin gắn với một người cụ thể hoặc giúp nhận ra người đó: họ tên, ngày sinh, số điện thoại, số căn cước, địa chỉ, và cả tài khoản mạng xã hội hay thư điện tử của người đó. Nghị định 13/2023/NĐ-CP xếp “thông tin về tài khoản số của cá nhân” vào nhóm dữ liệu cá nhân cơ bản (Điều 2 khoản 3 điểm k). Nghĩa là một tài khoản Facebook hay Gmail, xét theo pháp luật, cũng là dữ liệu cá nhân của chủ tài khoản.

Với người làm quảng cáo hay kẻ lừa đảo, những thông tin này có giá. Có danh sách số điện thoại kèm họ tên, họ có thể gọi đúng tên nạn nhân để tạo lòng tin. Có tài khoản mạng xã hội đã “nuôi” lâu năm, họ có thể chạy quảng cáo, đăng tin giả hoặc nhắn tin lừa đảo mà ít bị nền tảng chặn. Vì có người mua, dữ liệu cá nhân dần được mua đi bán lại như một mặt hàng.

Việc này ở Việt Nam không mới. Năm 2021, Bộ Công an triệt phá một đường dây mua bán gần 1.300 GB dữ liệu, chứa hàng tỷ thông tin về cá nhân và tổ chức, được rao bán công khai trên Facebook, Zalo, Telegram và diễn đàn tin tặc~\cite{sggp-2021}. Tháng 7/2024, Cục An ninh mạng và phòng, chống tội phạm sử dụng công nghệ cao (A05) cho biết đã phát hiện hàng trăm cá nhân, tổ chức liên quan đến việc bán dữ liệu cá nhân, với khối lượng hàng nghìn GB. Dữ liệu được bán công khai, kéo dài, có cả cam kết “bảo hành” và cập nhật theo yêu cầu người mua~\cite{vneconomy-2024}. Đến tháng 7/2026, Công an Nghệ An vẫn bắt giữ một nhóm mua bán hơn 50 triệu thông tin cá nhân qua Facebook và các nhóm kín trên Telegram~\cite{sggp-2026-50trieu}.

Hậu quả rõ nhất của việc mua bán dữ liệu là lừa đảo trực tuyến. Theo Hiệp hội An ninh mạng quốc gia, trong 11 tháng đầu năm 2025, thiệt hại do lừa đảo trực tuyến ước tính trên 6.000 tỷ đồng; 88,05\% người dùng được khảo sát cho biết từng bị mời chào dịch vụ mà họ không đăng ký, một dấu hiệu cho thấy thông tin của họ đã bị lộ~\cite{sggp-2026-thiethai}.

Các vụ kể trên phần lớn là nhóm nhỏ, bán qua hội nhóm kín. Vụ taphoammo.net khác ở chỗ việc mua bán được tổ chức thành một sàn giao dịch, vận hành như một doanh nghiệp.

\subsubsection{Vụ taphoammo.net}\label{sec:1.1.2}

“MMO” là viết tắt của \textit{make money online}, cách gọi chung các hoạt động kiếm tiền trên mạng. Trong giới này, “tài nguyên MMO” thường là tài khoản mạng xã hội, tài khoản thư điện tử, dữ liệu khách hàng và các công cụ phục vụ chạy quảng cáo hay tạo tài khoản hàng loạt. taphoammo.net (“Tạp hóa MMO”) là nơi mua bán những thứ đó.

Theo thông báo của Công an TP Hà Nội ngày 9 và 10/7/2026, website này do Đào Quang Huy (sinh năm 1992), Giám đốc Công ty TNHH Phát triển Công nghệ số HCT, cùng đồng phạm lập và vận hành từ cuối năm 2021~\cite{cand-2026,vnn-2026}. Nhóm tự giới thiệu là một công ty khởi nghiệp công nghệ, có đủ các bộ phận quản trị, kỹ thuật, marketing, chăm sóc khách hàng, kế toán và thanh toán.

Sàn hoạt động theo mô hình trung gian, gần giống một sàn thương mại điện tử thông thường. Người bán mở gian hàng để rao bán dữ liệu cá nhân, tài khoản mạng xã hội, thư điện tử, dịch vụ xác thực tài khoản và các tài nguyên số khác. Người mua chọn hàng và trả tiền qua sàn. Theo cơ quan điều tra, sàn kiểm duyệt người bán, bảo đảm giao dịch, giải quyết tranh chấp giữa các bên, và thu phí trung gian 3 đến 5\% trên mỗi giao dịch~\cite{vnn-2026}. Nói cách khác, sàn không chỉ “cho thuê chỗ” mà trực tiếp tham gia vào từng giao dịch.

Để tránh bị phát hiện, nhóm thuê tên miền và đặt máy chủ ở nước ngoài. Sàn được quảng bá trên Facebook, Telegram, TikTok, Instagram và Google.

Quy mô theo số liệu điều tra ban đầu:

{\small\setstretch{1.1}
\begin{table}[H]
\centering
\caption{Quy mô sàn taphoammo.net theo số liệu điều tra ban đầu}\label{tab:1.1}
\begin{tabularx}{\textwidth}{>{\hsize=0.945\hsize\raggedright\arraybackslash}X>{\hsize=1.055\hsize\raggedright\arraybackslash}X}
\toprule
\textbf{Chỉ số} & \textbf{Số liệu sơ bộ} \\
\midrule
Tài khoản đăng ký & Hơn 1,35 triệu \\
Gian hàng & Hơn 46.000 (trên 3.000 đang hoạt động) \\
Giao dịch & Hơn 53 triệu \\
Doanh thu (tổng tiền giao dịch qua sàn) & Trên 400 tỷ đồng \\
Phí trung gian & 3 đến 5\% mỗi giao dịch \\
Phí trung gian nhóm điều hành giữ lại & Nhiều tỷ đồng \\
\bottomrule
\end{tabularx}
\end{table}}

Nguồn:~\cite{cand-2026,vnn-2026,hnm-2026}.

Theo cơ quan điều tra, dữ liệu và tài khoản mua trên sàn đã bị dùng để lừa đảo chiếm đoạt tài sản, giả mạo danh tính và tạo tài khoản ảo. Đã có trường hợp đối tượng lừa đảo công nghệ cao dùng tài khoản mua từ hệ thống này~\cite{cand-2026}.

\subsubsection{Diễn biến vụ việc}\label{sec:1.1.3}

{\small\setstretch{1.1}
\begin{table}[H]
\centering
\caption{Diễn biến vụ việc}\label{tab:1.2}
\begin{tabularx}{\textwidth}{>{\raggedright\arraybackslash}l>{\hsize=1.000\hsize\raggedright\arraybackslash}X}
\toprule
\textbf{Thời điểm} & \textbf{Sự kiện} \\
\midrule
Cuối năm 2021 & Nhóm lập và bắt đầu vận hành taphoammo.net \\
6/2025 & Phòng An ninh mạng và phòng, chống tội phạm sử dụng công nghệ cao (Công an TP Hà Nội) phát hiện qua công tác nắm tình hình trên không gian mạng, xác lập chuyên án \\
25/12/2025 & Cơ quan Cảnh sát điều tra Công an TP Hà Nội ra quyết định khởi tố vụ án hình sự \\
2/2026 & 6 tổ công tác đồng loạt khám xét, triệu tập đối tượng tại Hà Nội, Phú Thọ, Quảng Ninh; thu giữ 20 máy tính, 25 điện thoại, nhiều ô tô cùng tài liệu, dữ liệu điện tử tại trụ sở công ty ở khu đô thị The Manor Nguyễn Xiển (Hà Nội) \\
9–10/7/2026 & Công an TP Hà Nội công bố khởi tố vụ án và 3 bị can về tội “Đưa hoặc sử dụng trái phép thông tin mạng máy tính, mạng viễn thông” (Điều 288 Bộ luật Hình sự) \\
\bottomrule
\end{tabularx}
\end{table}}

Nguồn:~\cite{cand-2026,k14-2026,hnm-2026}.

Từ khi lập chuyên án đến khi khám xét mất khoảng 8 tháng. Cơ quan điều tra nêu rõ việc hạ tầng đặt ở nước ngoài gây khó khăn cho việc truy vết.

Về mốc khởi tố, thông tin của Bộ Công an~\cite{bca-2026} và Báo Pháp luật Việt Nam~\cite{bpl-2026} (cùng ngày 14/7/2026) ghi Cơ quan Cảnh sát điều tra Công an TP Hà Nội ra quyết định khởi tố vụ án hình sự ngày 25/12/2025.

\subsubsection{Tình trạng tố tụng}\label{sec:1.1.4}

Theo công bố tháng 7/2026, vụ án đang được điều tra mở rộng. Chưa có kết luận điều tra, cáo trạng hay bản án. Đến ngày 22/10/2026, nhóm chưa tìm thấy thông tin công khai nào mới hơn. Các số liệu nêu trên đều là số liệu sơ bộ và có thể thay đổi khi có kết luận chính thức.

\subsubsection{Lý do chọn tình huống}\label{sec:1.1.5}

Nhóm chọn vụ này vì bốn lý do sau:

\begin{enumerate}
\item \textbf{Liên quan trực tiếp đến Luật An ninh mạng 2025.} Đề tài của nhóm xoay quanh Luật An ninh mạng 2025 (số 116/2025/QH15)~\cite{luat-anm-2025}, và theo nhận định của nhóm, nhiều hành vi của sàn taphoammo.net có dấu hiệu thuộc nhóm hành vi mà luật này cấm. Cụ thể, luật cấm kinh doanh trái phép dữ liệu cá nhân của người khác (điểm h khoản 2 Điều 7), cấm mua bán trái phép thông tin tài khoản số (điểm g khoản 4 Điều 13) và cấm lừa đảo chiếm đoạt tài sản (điểm đ khoản 2 Điều 7). Vì hành vi xảy ra trước 1/7/2026 nên vụ án vẫn xử theo luật cũ (phân tích ở Chương 3). Tuy vậy, nhóm muốn dùng vụ này để chỉ ra nếu bây giờ có một sàn tương tự thì sẽ vi phạm những điều nào của luật mới (Chương 7).
\item \textbf{Rút ra được cách tuân thủ.} Luật An ninh mạng 2025 yêu cầu doanh nghiệp cung cấp dịch vụ trên mạng phải có biện pháp ngăn chặn và kịp thời gỡ bỏ thông tin xâm phạm bí mật cá nhân, đời sống riêng tư (điểm b khoản 2 Điều 15), cung cấp thông tin và gỡ bỏ thông tin vi phạm theo yêu cầu của lực lượng chuyên trách (điểm a, b khoản 2 Điều 25) và lưu trữ dữ liệu người dùng tại Việt Nam (khoản 3 Điều 25). Sàn này thì ngược lại: có kiểm duyệt người bán nhưng vẫn cho rao bán dữ liệu cá nhân, còn máy chủ thì đặt ở nước ngoài. Từ những điểm sai này, nhóm có cơ sở để đề xuất cách một nền tảng nên làm để không vi phạm luật (Chương 5).
\item \textbf{Quy mô lớn và đủ thông tin để phân tích.} Đây không phải một nhóm nhỏ bán dữ liệu trong hội kín, mà là một sàn có hơn 1 triệu tài khoản, quảng cáo trên các nền tảng lớn và đứng tên một công ty có đăng ký. Cơ quan điều tra cũng đã công bố khá chi tiết về mô hình vận hành, quy mô và hạ tầng, đủ để đối chiếu từng hành vi với quy định.
\item \textbf{Kéo dài qua giai đoạn pháp luật thay đổi.} Từ cuối 2021 đến 2/2026, khung pháp luật về dữ liệu cá nhân và an ninh mạng thay đổi nhiều lần: Luật An toàn thông tin mạng 2015 và Luật An ninh mạng 2018, rồi Nghị định 13/2023, rồi Luật Bảo vệ dữ liệu cá nhân 2025. Đến 1/7/2026 thì Luật An ninh mạng 2025 có hiệu lực, thay cho hai luật về mạng trước đó. Vụ việc cho phép kiểm tra xem mỗi văn bản xử lý được hành vi tới đâu.
\end{enumerate}
\subsection{Phạm vi và cách tiếp cận của đề tài}\label{sec:1.2}

Một vụ việc lớn như taphoammo.net liên quan đến rất nhiều người và nhiều lĩnh vực pháp luật. Để phân tích có chiều sâu, nhóm giới hạn đề tài theo năm điểm dưới đây.

\textbf{Đối tượng phân tích.} Bài tập trung vào nền tảng taphoammo.net và nhóm vận hành nó (công ty HCT và các bị can). Hàng chục nghìn người bán, người mua trên sàn, cũng như những người dùng dữ liệu mua được để lừa đảo, chỉ được nhắc đến khi cần giải thích hành vi của sàn.

\textbf{Khung thời gian.} Từ cuối năm 2021, khi sàn bắt đầu hoạt động, đến tháng 2/2026, khi cơ quan công an khám xét. Đây là khoảng thời gian có hành vi được xác định. Các mốc tố tụng (quyết định khởi tố, công bố vụ án) được dùng làm thông tin về tình trạng tố tụng, không phải đối tượng phân tích.

\textbf{Lĩnh vực pháp luật.} Trọng tâm là hai lĩnh vực: bảo vệ dữ liệu cá nhân và an ninh mạng (bao gồm cả quy định hình sự về sử dụng trái phép thông tin trên mạng). Ba lĩnh vực khác là thương mại điện tử, quảng cáo và thanh toán chỉ được dùng bổ trợ, để phân tích các hành vi vận hành sàn, quảng bá trên mạng xã hội và thu tiền qua sàn (các hành vi về kiểm duyệt người bán, quảng cáo và thanh toán ở Chương 3). Danh sách văn bản và điều khoản cụ thể nằm ở mục~\ref{sec:2.2}.

\textbf{Cách tiếp cận.} Đề tài là một nghiên cứu tình huống (case study). Nhóm thu thập thông tin đã được công bố, chủ yếu từ thông báo của Công an TP Hà Nội được các báo chính thống đăng lại, và đối chiếu chéo giữa các nguồn. Dữ kiện nào chỉ xuất hiện ở một nguồn, hoặc các nguồn ghi khác nhau, được ghi chú rõ. Sau đó, từng hành vi được đặt cạnh quy định pháp luật có hiệu lực tại thời điểm hành vi xảy ra để xem có vi phạm hay không.

\textbf{Giới hạn.} Vụ án chưa có kết luận điều tra và chưa xét xử. Theo nguyên tắc suy đoán vô tội, chưa ai bị coi là có tội khi chưa có bản án có hiệu lực. Vì vậy, bài tách bạch hai loại nội dung:

\begin{itemize}
\item \textbf{Nội dung đã được cơ quan điều tra xác định}: các dữ kiện, con số và tội danh đã công bố chính thức.
\item \textbf{Nhận định của nhóm}: việc gắn từng hành vi với từng điều khoản cụ thể, và các hành vi có dấu hiệu vi phạm ngoài nội dung khởi tố.
\end{itemize}
Ngoài ra, nhóm không có quyền truy cập hồ sơ vụ án, nên một số thông tin chưa thể kiểm chứng, ví dụ sàn có đăng ký với Bộ Công Thương hay có giấy phép trung gian thanh toán hay không. Những điểm này được nêu rõ tại chỗ khi phân tích.

